\section{Function Call Mechanics and Parameter Binding}

A function definition creates a callable object, but the real execution machinery begins when that object is called. A call is not only a jump to another block of code. It is a structured operation: argument expressions are evaluated, object references are collected, a new frame is created, and parameter names inside that frame are bound according to the function signature.

This section studies that entrance mechanism in detail. First, it explains what is actually passed across the call boundary: not C-style variable boxes, not copied Python objects, but references to already existing objects bound to new local names. This is why reassignment inside a function stays local, while mutation of a shared mutable object can be visible to the caller.

After that, the section turns to signatures: positional arguments, keyword arguments, default values, positional-only parameters, keyword-only parameters, and annotations. These features are not decorative syntax. They are the rules by which Python constructs the initial local namespace of a function call.